\section{Lecture 6, Part 0: Bi-Encoders \& Contrastive Training}

\begin{frame}
\vfill
\centering
{\color{RLBlue}\large Lecture 6}

\vspace{0.3em}
{\color{RLBlue}\Large\bfseries Part 0}

\vspace{0.6em}
{\Large Bi-Encoders \& Contrastive Training}
\vfill
\end{frame}

\begin{frame}{Learning Goals}
\begin{itemize}
  \item Score a document with BM25 by hand, and say what $k_1$ and $b$ each control (Robertson \& Zaragoza, 2009).
  \pause
  \item Explain why a bi-encoder is the only design that scales to a whole corpus, and what it gives up against a cross-encoder.
  \pause
  \item Describe how in-batch and hard negatives train an embedding model (Karpukhin et al., 2020).
\end{itemize}
\end{frame}

% ─── BM25 first: the baseline that must be beaten ───────────────────────────
\begin{frame}{Why Start With Keyword Search?}
\begin{itemize}
  \item \textbf{Why:} the literature moved to dense retrieval, so it is tempting to skip the sparse baseline. Do not. On statute text a keyword index is fast, needs no training, no GPU and no model download -- and it is often hard to beat.
  \pause
  \item Legal drafting is \emph{formulaic}. An article about working hours really does contain the words ``working time'' and ``hours per week'', because the drafter was required to be unambiguous.
  \pause
  \item So BM25 is not the straw man in this lecture. It is the number your embedding model has to beat before it has earned its place in the pipeline.
  \pause
  \item \textbf{And it is the honest starting point:} you can build it in an afternoon, measure it, and only then find out whether the hard part is worth doing.
\end{itemize}
\end{frame}

\begin{frame}{BM25}
\begin{rldefinition}
\textbf{BM25} scores a document against a query by adding up, for each query term, how often it appears in \emph{this} document and how rare it is across the corpus:
\vspace{-0.3em}
\[
  \mathrm{score}(D,Q) \;=\; \sum_{q \in Q} \underbrace{\mathrm{IDF}(q)}_{\text{how rare}}
  \cdot \underbrace{\frac{f(q,D)\,(k_1+1)}{f(q,D) + k_1\!\left(1 - b + b\frac{|D|}{\mathrm{avgdl}}\right)}}_{\text{how often, with diminishing returns}}
\]
\vspace{-0.6em}

\textbf{$f(q,D)$} -- times the term occurs here. \emph{``week'' occurs once in Article 182.}\\
\textbf{IDF} -- rarity, $\ln\!\big(1 + \frac{N - n(q) + 0.5}{n(q)+0.5}\big)$. \emph{``week'' is in 1 of 5 articles; ``working'' in 3.}\\
\textbf{$k_1$} -- how fast repeats stop helping, typically $1.2$--$2.0$.\\
\textbf{$b$} -- how much to punish long documents, typically $0.75$.
\end{rldefinition}
\end{frame}

\begin{frame}{Worked Example: Ranking Five Articles}
\small
\textbf{Query:} \texttt{working time hours per week}. Five provisions, average length $15.2$ tokens.
\smallskip
\pause

\centering
\begin{tabular}{l|r|l}
Article & BM25 & what it is about \\
\hline
\textbf{182} & \textbf{5.85} & normal working time, 40 hours per week \\
183          & 1.46          & reduced working time for some categories \\
241          & 0.54          & annual leave \\
220          & 0.00          & payment of wages \\
99           & 0.00          & form of the employment contract \\
\end{tabular}

\vspace{0.6em}
\raggedright
\pause
\begin{itemize}
  \item The ranking is driven by \textbf{rarity}, not by counts: $\mathrm{IDF}(\text{``week''}) = 1.386$ against $\mathrm{IDF}(\text{``working''}) = 0.539$, because ``working'' appears in three of the five articles and ``week'' in one.
  \pause
  \item Two articles score exactly $0$ -- they share no query term at all. A sparse index does not rank them badly; it never sees them.
\end{itemize}
\end{frame}

\begin{frame}{Why Does BM25 Damp Repeated Terms?}
\begin{columns}[T]
\column{0.50\textwidth}
\small
\begin{itemize}
  \item \textbf{Why:} an article that says ``working time'' ten times is not ten times more about working time. Raw counts would let one long, repetitive document dominate every query.
  \pause
  \item The $k_1$ term saturates: at $k_1 = 1.2$ the first occurrence buys $0.46$ of the maximum, the second takes it to $0.63$, and the tenth only to $0.89$.
  \pause
  \item $b$ does the matching job for length -- dividing by $|D|/\mathrm{avgdl}$ so a long article cannot win on sheer size.
\end{itemize}
\column{0.48\textwidth}
\centering
\includegraphics[width=\linewidth]{figures/retrieval/retrieval_bm25_saturation.pdf}
\end{columns}
\end{frame}

\begin{frame}{Real Event: The Query That Returns Nothing}
\begin{columns}[T]
\column{0.50\textwidth}
\small
\begin{itemize}
  \item Same corpus, two queries with the \emph{same meaning}.
  \pause
  \item \texttt{working time hours week} $\rightarrow$ Article 182 at $4.53$, comfortably first.
  \pause
  \item \texttt{how long is the standard workweek} $\rightarrow$ after stopwords are removed the content words are \texttt{standard} and \texttt{workweek}. Neither appears anywhere in the corpus, so \textbf{every article scores exactly $0$}.
  \pause
  \item Article 182 is not ranked badly here. It is \emph{unreachable}: no threshold, no reranker and no amount of $k$ will surface a document the index scored zero.
\end{itemize}
\column{0.48\textwidth}
\centering
\includegraphics[width=\linewidth]{figures/retrieval/retrieval_lexical_gap.pdf}
\end{columns}
\end{frame}

\begin{frame}{The Apostrophe That Breaks Your Index}
\begin{rlpitfall}
Uzbek Latin writes \emph{o\textquoteleft} and \emph{g\textquoteleft} with \textbf{U+02BB}, and the tutuq belgisi with \textbf{U+02BC}. Both are Unicode category \texttt{Lm} -- \emph{letters}. The ASCII \texttt{\char39} (U+0027) is category \texttt{Po} -- punctuation. Three marks that look alike, three different codepoints -- which is why this failure is silent.
\end{rlpitfall}
\pause
\small
\centering
\begin{tabular}{l|l|c}
the word, written as & \texttt{\textbackslash w+} splits it into & matches? \\
\hline
\texttt{o\textquoteleft rtacha} \;(U+02BB, the corpus) & one token \;\texttt{o\textquoteleft rtacha} & --- \\
\texttt{o\char39 rtacha} \;(U+0027, what users type)          & two tokens \;\texttt{o} \;\texttt{rtacha} & \textcolor{red!65!black}{\textbf{no}} \\
\texttt{o\textquoteright rtacha} \;(U+2019, pasted from Word) & two tokens \;\texttt{o} \;\texttt{rtacha} & \textcolor{red!65!black}{\textbf{no}} \\
\end{tabular}

\vspace{0.5em}
\raggedright
\pause
\begin{itemize}
  \item \textbf{Normalisation does not fix this.} NFC, NFD, NFKC and NFKD all leave U+02BB exactly as it was -- it is a letter in its own right, not a decorated form of anything.
  \pause
  \item The bug is \textbf{normalising the corpus and the query differently}. Fold both sides with one rule, and test that a keyboard-typed query still finds its article.
\end{itemize}
\end{frame}

% ─── Bi-encoders ────────────────────────────────────────────────────────────
\begin{frame}{Why Do We Need a Bi-Encoder?}
\begin{itemize}
  \item \textbf{Why:} the obvious fix for the lexical gap is to let a model read the query and the article \emph{together} and judge the match. That is a \textbf{cross-encoder}, and it is the most accurate thing we know how to build.
  \pause
  \item It is also unusable as a first stage. Judging a pair requires one forward pass, and the pair changes with every query -- so nothing can be precomputed.
  \pause
  \item For the capstone corpus that is \textbf{7{,}368 model passes for one question}. At the other end of the room, a user is waiting.
  \pause
  \item A \textbf{bi-encoder} gives that up deliberately: encode the query and the article \emph{separately}, so every article can be encoded once, offline, and a query becomes one pass plus a pile of dot products.
\end{itemize}
\end{frame}

\begin{frame}{Bi-Encoder vs.\ Cross-Encoder}
\centering
\begin{tikzpicture}[
  font=\scriptsize,
  enc/.style={draw=RLBlue, fill=RLGray, rounded corners=2pt, inner sep=3.5pt,
              minimum height=5.5mm, minimum width=17mm, align=center},
  txt/.style={draw=RLGreen, fill=RLGray, rounded corners=2pt, inner sep=3pt,
              minimum height=5mm, minimum width=15mm, align=center},
  res/.style={draw=RLOrange, fill=RLGray, rounded corners=2pt, inner sep=3pt,
              minimum height=5mm, minimum width=15mm, align=center},
  lbl/.style={font=\scriptsize, inner sep=1.5pt}
]
  % --- bi-encoder ---
  \node[lbl, anchor=west] at (-0.4,2.9) {\textbf{Bi-encoder} -- encoded apart, compared by a dot product};
  \node[txt] (q1) at (0.6,2.0)  {query};
  \node[txt] (d1) at (0.6,0.85) {article};
  \node[enc] (e1) at (3.3,2.0)  {encoder};
  \node[enc] (e2) at (3.3,0.85) {encoder};
  \node[lbl] (v1) at (5.7,2.0)  {$\mathbf{q}\in\mathbb{R}^{d}$};
  \node[lbl] (v2) at (5.7,0.85) {$\mathbf{a}\in\mathbb{R}^{d}$};
  \node[res] (s1) at (8.3,1.42) {$\mathbf{q}\cdot\mathbf{a}$};
  \draw[->] (q1) -- (e1);  \draw[->] (d1) -- (e2);
  \draw[->] (e1) -- (v1);  \draw[->] (e2) -- (v2);
  \draw[->] (v1) -- (s1);  \draw[->] (v2) -- (s1);
  \node[lbl, anchor=west, text=RLGreen] at (0.0,0.15)
    {the article side runs \emph{once}, offline -- nothing on it depends on the query};

  % --- cross-encoder ---
  \node[lbl, anchor=west] at (-0.4,-0.9) {\textbf{Cross-encoder} -- read together, one pass per pair};
  \node[txt] (p) at (1.4,-1.9) {query $\oplus$ article};
  \node[enc] (e3) at (4.6,-1.9) {encoder};
  \node[res] (s2) at (7.6,-1.9) {score};
  \draw[->] (p) -- (e3);   \draw[->] (e3) -- (s2);
  \node[lbl, anchor=west, text=red!65!black] at (8.6,-1.9) {$\times N$};
\end{tikzpicture}

\vspace{0.3em}
\footnotesize Everything else in this lecture follows from that gap: the corpus is encoded once and kept, so a query costs one model pass plus arithmetic.
\end{frame}

\begin{frame}{Embedding, Bi-Encoder, Cosine Similarity}
\begin{rldefinition}
\textbf{Embedding} -- a fixed-length vector standing for a piece of text, built so that nearness in the vector space means nearness in meaning. \emph{Here:} Article 182 becomes one vector of $d$ numbers.

\smallskip
\textbf{Bi-encoder} -- two encoder passes that never see each other's input, scored by a dot product. \emph{Here:} one tower for questions, one for articles.

\smallskip
\textbf{Cosine similarity} -- the dot product of unit-length vectors, $\cos\theta = \frac{\mathbf{q}\cdot\mathbf{a}}{\lVert\mathbf{q}\rVert\lVert\mathbf{a}\rVert}$: the angle between them, ignoring length. \emph{Why:} a long article should not outscore a short one for being long.
\end{rldefinition}
\pause
\begin{itemize}
  \item Normalise once at index time and cosine \emph{is} the dot product -- which is why libraries ask you to store unit vectors.
\end{itemize}
\end{frame}

\begin{frame}{The Cost of Each Design}
\begin{columns}[T]
\column{0.50\textwidth}
\small
\begin{itemize}
  \item \textbf{Bi-encoder:} $N$ encoder passes \emph{once, offline}, then per query one pass plus $N$ dot products. A dot product is arithmetic, not a model.
  \pause
  \item \textbf{Cross-encoder:} $N$ encoder passes \emph{per query}, every query, forever.
  \pause
  \item At $N = 7{,}368$ that is one pass against $7{,}368$. At a million articles it is one against a million.
  \pause
  \item \textbf{So use both.} A bi-encoder retrieves the top few dozen, a cross-encoder reorders only those -- the reranking stage of Lecture 5. Accuracy where it is affordable.
\end{itemize}
\column{0.48\textwidth}
\centering
\includegraphics[width=\linewidth]{figures/retrieval/retrieval_encoder_cost.pdf}
\end{columns}
\end{frame}

% ─── Contrastive training ───────────────────────────────────────────────────
\begin{frame}{Why Does Contrastive Training Work?}
\begin{itemize}
  \item \textbf{Why:} an untrained encoder has no reason to put a \emph{question} near the \emph{article that answers it} -- they do not look alike. That correspondence has to be taught.
  \pause
  \item \textbf{The training signal:} take pairs (question, article that answers it). Pull each pair together, push everything else apart. Nothing else is needed -- no relevance grades, no ranking labels.
  \pause
  \item \textbf{In-batch negatives} make ``everything else'' free. In a batch of $64$ pairs, each question already has its own article plus $63$ others sitting in the same batch. Score all $64$, apply cross-entropy, and the correct one must win.
  \pause
  \item So a batch of $64$ costs $128$ encoder passes and yields $64$ classification problems over $64$ candidates each. \textbf{Bigger batches are better training} here -- more negatives per step, at no extra encoding cost.
\end{itemize}
\end{frame}

\begin{frame}{Why Hard Negatives Matter}
\begin{itemize}
  \item \textbf{The flaw in free negatives:} a random other article in the batch is about annual leave or contract form. Telling them apart from a working-hours question is trivial, so after a while the gradient carries almost nothing.
  \pause
  \item \textbf{Hard negatives} are articles that \emph{look} right and are not -- fetch the BM25 top results for the question and take the ones that are not the gold article. Article 183, ``reduced working time'', is exactly this for our query.
  \pause
  \item Karpukhin et al.\ (2020) train DPR this way -- gold pair, one BM25 hard negative, plus in-batch negatives -- and report beating BM25 by $9$--$19$ points of top-$20$ retrieval accuracy across open-domain QA benchmarks.
  \pause
  \item \textbf{Note what that implies for you:} the sparse index you built first is not thrown away when you train a dense model. It is what \emph{generates the training data}.
\end{itemize}
\end{frame}
